\documentclass[a4paper,11pt]{article}
\usepackage{sbpo-template}
\usepackage[brazil]{babel}    % português do Brasil
\usepackage[utf8]{inputenc}   % acentos em UTF-8
\usepackage[T1]{fontenc}      % para hifenização correta
\usepackage{amsmath,amssymb}
\usepackage{url}
\usepackage[square]{natbib}
\usepackage{indentfirst}
\usepackage{fancyhdr}
\usepackage{graphicx}
\usepackage{booktabs}
\usepackage{makecell}
\pagestyle{fancy}
\fancyhf{}
\fancyhead[C]{
\includegraphics[scale=0.5]{figuras/sbpo2026-header-logo.png}}
%\renewcommand{\headrulewidth}{0pt}
\renewcommand{\headruleskip}{-1mm}
%\setlength\headheight{101.0pt}
%\addtolength{\textheight}{-101.0pt}
\setlength\headheight{86pt}
\addtolength{\textheight}{-86pt}
\setlength{\headsep}{5mm}
\setlength{\footskip}{4.08003pt}

\title{Otimização de Inventário Multi-Echelon sob Demanda Intermitente: Comparação de Doze Políticas em Dados Reais de Varejo Brasileiro}

\author{
\name{}
\institute{Filiação}
\iaddress{Endereço da Instituição}
\email{e-mail}
}

\author{
\name{Segundo Autor}
\institute{Filiação}
\iaddress{Endereço da Instituição}
\email{e-mail}
}

\begin{document}

\maketitle

\vspace{8mm}
\begin{resumo}
Modelos clássicos de inventário, como o \textit{Economic Order Quantity}
(EOQ) e a política de dois limiares~$(s,S)$, pressupõem
estacionariedade que falha sistematicamente em cenários com demanda
intermitente, gerando excesso de estoque e rupturas frequentes no
varejo regional brasileiro. Estudos comparativos recentes baseados em
conjuntos como o Walmart~M5 não cobrem o regime de alta variabilidade
caracterizado por coeficiente de variação (CV) entre $2.5$ e $5.4$,
típico de produtos de baixo giro. Este trabalho apresenta uma
comparação de doze políticas de controle de inventário em dados reais
de varejo brasileiro, avaliadas em cinco indicadores: custo total,
nível de serviço, taxa de ruptura, efeito de amplificação da demanda
(\textit{bullwhip}) e frequência de pedidos.
A arquitetura híbrida GA-DQN (\textit{Genetic Algorithm with Deep Q-Network}) obtém o melhor equilíbrio prático: custo
substancialmente inferior às heurísticas clássicas, efeito
\textit{bullwhip} significativamente menor que a política~$(s,S)$ e
taxa de ruptura equivalente à melhor política prática isolada,
confirmando a robustez da abordagem em regime de extrema intermitência.
\end{resumo}


\bigskip
\begin{palchaves}
Gestão de inventário; Demanda intermitente; Cadeia de suprimentos
multi-nível; Meta-heurísticas; Aprendizado por reforço.

\bigskip
\noindent{Tópicos (em ordem de PRIORIDADE): }
Logística e \textit{Supply Chain}; Otimização de Processos;
Inteligência Artificial e Aprendizado de Máquina.
\end{palchaves}

\vspace{8mm}
\begin{abstract}
Classical inventory models, such as the Economic Order Quantity (EOQ)
and the two-threshold policy~$(s,S)$, rely on stationarity assumptions
that systematically fail for intermittent-demand items, producing
excess stock and frequent stockouts in Brazilian regional retail.
Recent comparative studies based on datasets such as Walmart~M5 do
not cover this high-variability regime, characterized by a coefficient
of variation (CV) between $2.5$ and $5.4$, typical of slow-moving
products. This work presents a comparison of twelve inventory
control policies on real Brazilian retail data, evaluated on five
indicators: total cost, service level, stockout rate, demand
amplification along the supply chain (bullwhip effect), and order
frequency. The hybrid GA-DQN (Genetic Algorithm with Deep Q-Network)
architecture achieves the best practical trade-off: substantially
lower cost than classical heuristics, a bullwhip effect significantly
lower than $(s,S)$, and a stockout rate on par with the best practical
standalone policy, confirming the robustness of the hybrid approach
under extreme intermittent demand.
\end{abstract}

\bigskip
\begin{keywords}
Inventory management; Intermittent demand; Multi-level supply chain;
Metaheuristics; Reinforcement learning.

\bigskip
\noindent{Paper topics (in order of PRIORITY): }
Logistics and Supply Chain; Process Optimisation;
Artificial Intelligence and Machine Learning.
\end{keywords}

% =========================================================
\section{Introdução}
% =========================================================

A gestão ineficiente de inventário representa, segundo estimativas,
entre \(3.5\%\) e \(8\%\) do faturamento anual do varejo brasileiro em
capital imobilizado e perdas de venda por ruptura~\citep{silver:2017}.
Em redes com demanda intermitente, esse custo é ainda maior: itens
de baixo giro com coeficiente de variação (CV) elevado são
gerenciados, na maioria das redes regionais brasileiras, com as
mesmas heurísticas formuladas há mais de cem anos: o Lote Econômico
de Compra (EOQ)~\citep{harris:1913} e a política de dois
limites~$(s,S)$~\citep{hadley:1963}. Essas abordagens pressupõem
estacionariedade que raramente se verifica em dados reais de
varejo~\citep{silver:2017}, gerando excesso de estoque imobilizado,
rupturas frequentes e amplificação da demanda ao longo da cadeia.

O desafio acentua-se quando o padrão de consumo é
\emph{intermitente}: longos períodos de demanda nula alternados com
ocorrências esporádicas de alta magnitude~\citep{croston:1972}.
Esse comportamento é quantificado pelo coeficiente de variação
$CV = \hat{\sigma}_D / \hat{\mu}_D$, onde $\hat{\sigma}_D$ é o
desvio padrão e $\hat{\mu}_D$ é a média da demanda; valores acima de \(0.5\)
indicam demanda errática e acima de \(1.5\), demanda extremamente
intermitente~\citep{syntetos:2001}. Todos os 15 cenários deste
estudo possuem $CV$ entre \(2.5\) e \(5.4\), regime em que as hipóteses
clássicas falham de forma sistemática.

% Do ponto de vista gerencial, a escolha da política de inventário
% impacta diretamente três grupos de stakeholders: (i)~gestores de
% \emph{supply chain} e compradores de redes de varejo, que
% necessitam de critério objetivo para decidir quando abandonar
% heurísticas clássicas em favor de métodos adaptativos; (ii)~pequenas
% e médias empresas sem equipe dedicada de \emph{data science}, para
% quem o custo de implementação de uma solução complexa deve ser
% justificado por ganho mensurável; e (iii)~pesquisadores em Pesquisa
% Operacional e aprendizado de máquina, que demandam uma base de
% comparação reproduzível em regime de alta intermitência, complementar
% ao amplamente utilizado Walmart~M5.

Diante desse cenário, este trabalho ataca o problema \textit{Multi Echelon Inventory Optimization} (MEIO) com demanda intermitente: dado
um armazém regional que abastece 15 lojas varejistas ao longo de 38
ciclos bimestrais, determinar, para cada loja, a política de reposição
que minimiza o custo total acumulado (manutenção, ruptura e ordenação)
sob $CV$ entre $2.5$ e $5.4$. Para isso, apresenta-se a primeira
comparação sistemática de doze políticas em dados reais de varejo
brasileiro, regime não coberto por estudos baseados no Walmart~M5.
O estudo estende a referência mais próxima em cinco direções concretas,
descritas a seguir.

\citet{zabraoui:2025} constituem a referência mais próxima deste
estudo: os autores propuseram a arquitetura híbrida GA-DQN
(\textit{Genetic Algorithm-Deep Q-Network}) no
Walmart~M5, elevando o nível de serviço de $61\%$ para $94\%$ frente ao
DQN isolado. O presente artigo parte desse resultado e o estende
em cinco direções concretas:

\begin{itemize}
  \item Regime de maior intermitência: $CV$ entre $2.5$ e $5.4$,
    versus CV típico $< 1$ no Walmart~M5; os modelos de
    \citet{zabraoui:2025} não foram validados nesse regime.
  \item Escopo comparativo ampliado: incorporação de $(s,S)$,
    Jornaleiro, PSO (\textit{Particle Swarm Optimization}),
    Evolução Diferencial,
    SARSA (\textit{State-Action-Reward-State-Action})
    e a variante híbrida GA-PPO, ausentes em \citet{zabraoui:2025}.
  \item Avaliação primária em cinco indicadores: efeito
    \textit{bullwhip} e frequência de pedidos aparecem em
    \citet{zabraoui:2025} apenas para políticas selecionadas; aqui
    são elevados a critérios primários integrados na comparação
    sistemática de todas as doze políticas, incluindo taxa de ruptura.
  \item Estratificação por regime de demanda: três grupos de
    volume revelam como a política ótima depende do regime.
  \item Calibração adaptativa por loja: parâmetros
    escalonados individualmente para lojas com volume entre
    $11.7$ e $1455$ unidades por ciclo.
\end{itemize}

O restante do artigo está organizado como se segue: a Seção~2
revisa a literatura; a Seção~3 descreve os dados; a Seção~4 detalha
as doze políticas e seus hiperparâmetros; a Seção~5 apresenta a
metodologia experimental; a Seção~6 discute os resultados; e a
Seção~7 conclui com recomendações práticas.

% =========================================================
\section{Revisão da Literatura}
% =========================================================

Esta seção revisa os trabalhos mais relacionados, organizados em
duas frentes: (i)~aplicação de aprendizado por reforço (RL) ao controle
de inventário e (ii)~uso de meta-heurísticas sob demanda intermitente.
A revisão culmina na análise de \citet{zabraoui:2025}, referência
central deste trabalho; a fundamentação teórica de cada política é
detalhada na Seção~4.

\citet{oroojlooyjadid:2022} foram os primeiros a demonstrar, no
\textit{Beer Game} \textit{multi-echelon}, que a DQN aprende políticas de reposição
sem modelar explicitamente a distribuição de demanda, superando
políticas de \textit{base-stock} em cenários com informação assimétrica.
\citet{gijsbrechts:2022} ampliaram o escopo, comparando o
A3C (\textit{Asynchronous Advantage Actor-Critic}) contra
heurísticas de referência em três problemas canônicos (perdas de
venda, duplo abastecimento e \textit{multi-echelon}), concluindo que o RL
profundo pode igualar, mas não domina universalmente, as melhores
heurísticas. Para o problema de reposição conjunta,
\citet{vanvuchelen:2020} aplicaram o PPO (\textit{Proximal Policy Optimization})
com sucesso, superando políticas clássicas em custo total. Mais recentemente,
\citet{liu:2024} estenderam o paradigma ao ambiente descentralizado
com HAPPO (\textit{Heterogeneous-Agent Proximal Policy Optimization})
multi-agente, reduzindo simultaneamente custo e efeito
\textit{bullwhip} em cadeias \textit{multi-echelon}. Uma limitação
transversal a todos esses estudos é a exigência de horizontes longos
para convergência (da ordem de $10^6$ interações), tornando-os
inadequados para séries históricas curtas como as deste trabalho.

Em paralelo, uma segunda frente investiga meta-heurísticas
especificamente sob demanda intermitente. \citet{yuna:2023}
compararam \textit{Simulated Annealing} e Busca Tabu na otimização
de parâmetros $(s,S)$ em sete conjuntos de dados intermitentes,
verificando superioridade do SA em cinco dos sete cenários.
\citet{erkayman:2025} modelaram a chegada de demanda como cadeia de
Markov e otimizaram os parâmetros de reposição com GA e PSO,
mostrando que o PSO converge mais rapidamente em espaços de baixa
dimensão. Ambos os estudos, porém, restringem-se a um único escalão
e não incorporam agentes de RL, deixando em aberto a comparação
integrada em ambiente \textit{multi-echelon} de alta intermitência.

\citet{zabraoui:2025} realizaram o estudo comparativo mais próximo do
presente trabalho. Sobre o conjunto Walmart~M5, os autores compararam
cinco famílias de abordagens: heurísticas clássicas, modelos de
aprendizado de máquina convencional, redes neurais profundas,
algoritmos evolutivos e agentes de aprendizado por reforço. A
principal contribuição foi a arquitetura híbrida GA-DQN: o
Algoritmo Genético otimiza globalmente os parâmetros estáticos de
inventário (ponto de reposição e estoque de segurança), enquanto a
Rede Q Profunda aprende políticas adaptativas de pedido dependentes
do estado, refinando localmente as decisões a cada ciclo. Essa
combinação elevou o nível de serviço de $61\%$ (DQN isolado) para $94\%$,
com redução simultânea do custo total. As técnicas empregadas derivam
de quatro vertentes canônicas: teoria clássica de controle de
estoques~\citep{harris:1913,hadley:1963,cachon:2009}, computação
evolutiva~\citep{holland:1975}, aprendizado por reforço
profundo~\citep{mnih:2015} e controle de inventário por
RL~\citep{oroojlooyjadid:2022}. O presente trabalho parte dessas
mesmas bases, porém supera três limitações: (i) o Walmart~M5
apresenta CV típico $< 1$, regime não representativo do varejo
regional brasileiro ($CV$ entre $2.5$ e $5.4$ neste estudo);
(ii) o portfólio de \citet{zabraoui:2025} exclui PSO, Evolução
Diferencial, SARSA e a variante GA-PPO; e (iii) o efeito \textit{bullwhip} e a frequência de pedidos, embora
presentes em \citet{zabraoui:2025} para políticas selecionadas, não
constituem critérios primários: \citet{zabraoui:2025} concentra-se
em custo e nível de serviço. Este trabalho eleva esses indicadores a
critérios primários e acrescenta a taxa de ruptura, inédita nesse
escopo comparativo.

A Tabela~\ref{tab:relwork} sintetiza as dimensões em que este estudo
avança o estado da arte.

\begin{table}[ht]
\centering
\small
\caption{Comparação de trabalhos relacionados em gestão de inventário.
\textbf{Negrito}: dimensões em que este trabalho avança o estado da arte.
``Dados Reais'': dados proprietários coletados no domínio-alvo
(varejo com alta intermitência); Walmart~M5, embora real, não
representa esse regime (CV típico $<1$).}
\label{tab:relwork}
\begin{tabular}{@{}lccccc@{}}
\toprule
Trabalho &
\makecell[c]{Dados\\Reais} &
\makecell[c]{CV\\(regime)} &
\makecell[c]{Políticas\\(\#)} &
\makecell[c]{Multi-\\nível} &
\makecell[c]{Indicadores\\(\#)} \\
\midrule
\citealt{oroojlooyjadid:2022} & Não & $<1$         &  3 & Sim & 2 \\
\citealt{vanvuchelen:2020}    & Não & $<1$         &  4 & Não & 1 \\
\citealt{gijsbrechts:2022}    & Não & $<1$         &  6 & Sim & 2 \\
\citealt{yuna:2023}           & Não & $>1.5$     &  2 & Não & 2 \\
\citealt{erkayman:2025}       & Não & $>1.5$     &  2 & Não & 1 \\
\citealt{liu:2024}            & Não & $<1$         &  4 & Sim & 2 \\
\citealt{zabraoui:2025}       & Não & $<1$         &  7 & Sim & 3 \\
\midrule
\textbf{Este trabalho}
  & \textbf{Sim}
  & $\mathbf{2.5}$--$\mathbf{5.4}$
  & \textbf{12}
  & \textbf{Sim}
  & \textbf{5} \\
\bottomrule
\end{tabular}
\end{table}



% =========================================================
\section{Arquitetura Multi-Nível e Caracterização dos Dados}
% =========================================================

Esta seção descreve a cadeia de suprimentos modelada, o contexto e
a escala do conjunto de dados original, os critérios de seleção de
lojas e produto, o \textit{pipeline} de pré-processamento e a caracterização
estatística do regime de demanda das quinze unidades experimentais.
A formulação formal do problema é apresentada a seguir.

O problema \textit{Multi Echelon Inventory Optimization} (MEIO)
considera uma cadeia de suprimentos real com $N = 15$ lojas
varejistas abastecidas por um único armazém regional. O horizonte é
discretizado em ciclos bimestrais $t \in \{1, \ldots, T\}$, com
$T = 38$. As variáveis e parâmetros centrais são:

\begin{itemize}
  \item $D_{t,i} \geq 0$: demanda observada na loja $i$ no ciclo $t$;
  \item $I_{t,i}$: nível de estoque da loja $i$ no início de $t$;
  \item $Q_{t,i} \geq 0$: quantidade pedida em $t$, recebida após
    prazo $L = 2$ ciclos;
  \item $S_{t,i} = \max(0, D_{t,i} - I_{t,i})$: ruptura no ciclo $t$;
  \item $s_i$: ponto de reposição, acionado quando $I_{t,i} \leq s_i$;
  \item $h$, $c_s$, $c_o^{fix}$, $c_o^{var}$: custos unitários de
    manutenção, ruptura, ordenação fixa e variável.
\end{itemize}

A dinâmica de estoque é $I_{t+1,i} = \max(0, I_{t,i} - D_{t,i}) +
Q_{t-L,i}$. Uma \emph{política} $\pi$ mapeia o estado observado
$(I_{t,i},\, D_{t-1,i},\, Q_{t-L,i})$ em uma decisão de pedido
$Q_{t,i} = \pi(\cdot)$. O objetivo é determinar $\pi^*$ que minimize
o custo total acumulado:
\begin{equation}
  \min_\pi \sum_{t=1}^{T} \bigl[
    h\, I_{t,i} + c_s\, S_{t,i}
    + c_o^{fix}\,\mathbf{1}(Q_{t,i}>0)
    + c_o^{var}\, Q_{t,i}
  \bigr].
  \label{eq:obj}
\end{equation}

\subsection{Hierarquia de Três Níveis e Lógica de Decisão}

A cadeia de suprimentos modela um sistema de inventário multi-nível
com três escalões hierárquicos, estrutura referenciada na literatura
como \textit{Multi-Echelon Inventory Optimization}
(MEIO)~\citep{axsater:2006,clark:1960}. Os três escalões são:

\begin{itemize}
  \item Escalão 1 (Armazém Regional): único centro distribuidor da
    Paraíba; repõe junto ao fornecedor com prazo $L = 2$ ciclos.
  \item Escalão 2 (Lojas Revendedoras): 15 unidades varejistas com
    séries de demanda independentes; cada loja decide individualmente
    quando e quanto pedir ao armazém.
  \item Escalão 3 (Produto/SKU): produto 48130, único SKU estudado;
    foco em um item isola o comportamento intermitente de efeitos de
    portfólio e torna os 38 ciclos suficientes para calibração e avaliação.
\end{itemize}

As decisões de reposição de cada loja são tomadas de forma
descentralizada: cada política $\pi_i$ observa o estado local
$(I_{t,i}, D_{t-1,i}, Q_{t-L,i})$ e emite o pedido $Q_{t,i}$ sem
coordenação explícita com as demais lojas. Essa arquitetura de
decisão descentralizada é representativa da realidade operacional
da maioria das redes de varejo regional brasileiras, onde sistemas
integrados de planejamento de inventário multi-loja raramente
estão disponíveis.

\subsection{Contexto e Escala do Conjunto de Dados Original}

Os dados provêm de um sistema transacional de gestão de vendas de
uma rede de varejo do Nordeste brasileiro. O conjunto completo
contém aproximadamente 48 milhões de registros de transações de
venda, distribuídos entre cerca de 12 mil produtos únicos (SKUs) e
mais de 3 milhões de entidades de demanda (lojas revendedoras) em
múltiplos estados. Esses números evidenciam o regime de grande
escala em que o problema real se insere: o espaço de decisões cresce
combinatorialmente com o número de produtos, lojas e
ciclos~\citep{shapiro:2001,boute:2022}.

Para viabilizar o estudo experimental com rigor metodológico,
adotou-se um recorte operacional com três critérios de foco:
(i)~restrição geográfica ao estado da Paraíba (PB), que possui o
maior volume de lojas da rede nessa região; (ii)~restrição de
produto a um único SKU de comportamento representativo da demanda
intermitente; e (iii)~restrição temporal ao período disponível com
séries suficientemente longas. O recorte resulta em instâncias de
tamanho gerenciável que preservam a complexidade qualitativa do
problema de interesse.

\subsection{Estrutura Temporal e Ciclos de Venda}

O sistema registra vendas por \emph{ciclo bimestral} no formato
\texttt{AAAA-NC}, onde \texttt{AAAA} é o ano e \texttt{NC} é o
número do ciclo dentro do ano, de 1 a 17 (dezessete ciclos por
ano). Cada ciclo tem duração média de 21 dias corridos. A escolha da
granularidade bimestral reflete a política de reposição da rede:
pedidos ao armazém regional são consolidados por ciclo, o que
justifica o prazo de entrega $L = 2$ ciclos ($\approx 42$ dias)
adotado no modelo.

O horizonte temporal do experimento compreende os ciclos de
janeiro de 2023 a abril de 2025, totalizando $T = 38$ ciclos
consecutivos. Esse horizonte cobre dois anos completos mais quatro
meses, período suficiente para capturar a sazonalidade anual e os
padrões estruturais de intermitência de cada loja.

\subsection{Seleção do Produto e Critérios de Filtragem}

O produto 48130 foi selecionado a partir do conjunto completo de
12 mil SKUs com base em dois critérios ordenados: (i) maior
cobertura de lojas com demanda ativa, definida como o número de
lojas no estado com ao menos 17 ciclos de demanda positiva no
horizonte estudado; e (ii) regime de demanda extremamente
intermitente, caracterizado por CV $> 1.5$ em todas as lojas
ativas. O produto 48130 atende simultaneamente aos dois critérios,
com 33 lojas ativas no estado da Paraíba. Ele é representativo de
uma categoria de produto de baixo giro com padrão errático de
consumo, típico de artigos sazonais ou de reposição periódica em
pequenas quantidades.

\subsection{Pipeline de Pré-processamento}

O pré-processamento dos dados seguiu cinco etapas sequenciais:

\begin{itemize}
  \item Extração e filtragem geográfica: seleção de todos
    os registros com estado = PB e produto = 48130;
  \item Limpeza de anomalias: remoção de transações com
    quantidade negativa (devoluções registradas no mesmo campo de
    vendas) ou superior a 50.000 unidades (provável erro de entrada);
  \item Agregação por ciclo: soma das quantidades vendidas por
    (loja, ciclo), consolidando as transações diárias no
    nível de granularidade do modelo;
  \item Reconstituição de séries completas: para cada loja,
    criação de um vetor de comprimento $T = 38$ com imputação de zero
    nos ciclos sem registro, representando demanda nula nesses períodos;
  \item Seleção de lojas experimentais: ordenação das lojas
    pelo número de ciclos com demanda positiva (descendente) e
    seleção das quinze com ao menos 17 ciclos positivos.
\end{itemize}

O limiar mínimo de 17 ciclos positivos equivale a ao menos $45\%$
dos ciclos com demanda observada, condição necessária para que as
políticas baseadas em estatísticas históricas ($\hat{\mu}$,
$\hat{\sigma}$) tenham amostras suficientes para calibração.

\subsection{Caracterização Estatística da Demanda por Loja}

A Tabela~\ref{tab:demanda} apresenta o perfil completo das quinze
lojas selecionadas. Quatro estatísticas são reportadas: demanda média
apenas nos ciclos com demanda positiva $\hat{\mu}^+$ (excluindo
os zeros estruturais, para capturar a magnitude dos picos); desvio
padrão calculado sobre todos os 38 ciclos $\hat{\sigma}$ (incluindo
os zeros, capturando a variabilidade total); coeficiente de variação
$\text{CV} = \hat{\sigma}/\hat{\mu}$ (onde $\hat{\mu}$ é a média
sobre todos os 38 ciclos); e número de ciclos com demanda positiva.

Nenhuma loja apresenta $CV$ inferior a $2.5$, situando todos os
cenários bem acima do limiar de demanda extremamente intermitente
(CV $> 1.5$) de \citet{syntetos:2001}. Isso invalida as hipóteses
de normalidade subjacentes a EOQ e $(s,S)$: a distribuição empírica
de demanda é bimodal (zero com alta probabilidade; pico raro de
grande magnitude), e não gaussiana.

\begin{table}[ht]
\centering
\small
\caption{Perfil de demanda das quinze lojas selecionadas, produto
48130, Paraíba, 2023--2025. Lojas ordenadas por $\hat{\mu}^+$
decrescente. $\hat{\mu}^+$: média apenas nos ciclos com demanda
positiva; critérios de grupo baseados na média incondicional
$\hat{\mu}$ (sobre todos os 38 ciclos):
I ($\hat{\mu}>500$), II (50--500), III ($\leq 50$).}
\label{tab:demanda}
\begin{tabular}{@{}lrrrrc@{}}
\toprule
\makecell[l]{Loja} &
\makecell[r]{$\hat{\mu}^+$\\(unid./ciclo)} &
\makecell[r]{Desv. Pad.} &
\makecell[r]{CV} &
\makecell[r]{Ciclos\\Positivos} &
\makecell[c]{Grupo} \\
\midrule
1 & 1.455{,}2 & 5.037{,}5 & 5{,}14 & 18 & I \\
2  & 1.375{,}7 & 4.464{,}6 & 4{,}35 & 22 & I \\
3 &   703{,}6 & 2.561{,}1 & 5{,}39 & 18 & I \\
4 &   553{,}4 & 1.380{,}2 & 3{,}67 & 19 & I \\
5 &   518{,}2 & 1.150{,}1 & 2{,}83 & 25 & I \\
6  &   300{,}9 &   590{,}3 & 2{,}95 & 19 & II \\
7 &   246{,}0 &   546{,}4 & 3{,}04 & 22 & II \\
8 &   181{,}6 &   393{,}0 & 2{,}90 & 23 & II \\
9  &   169{,}7 &   489{,}5 & 3{,}79 & 23 & II \\
10 &    88{,}3 &   274{,}8 & 4{,}51 & 19 & II \\
11 &    85{,}4 &   253{,}3 & 4{,}31 & 19 & II \\
12 &    44{,}8 &    85{,}9 & 2{,}53 & 24 & III \\
13  &    43{,}1 &    88{,}4 & 2{,}83 & 22 & III \\
14 &    34{,}7 &    97{,}1 & 3{,}97 & 20 & III \\
15 &    11{,}7 &    25{,}5 & 3{,}24 & 19 & III \\
\bottomrule
\end{tabular}
\end{table}

A heterogeneidade entre lojas é expressiva: a demanda média varia
de $11.7$ (Loja~15) a $1455.2$ unidades por ciclo (Loja~1),
razão de 124 vezes. Esse gradiente inviabiliza
comparações com parâmetros absolutos uniformes e motivou a
calibração adaptativa descrita na Seção~5. A estratificação em três
grupos captura três regimes operacionalmente distintos:

\begin{itemize}
  \item Grupo I (5 lojas, $\hat{\mu} > 500$ unid./ciclo):
    lojas de alto volume com picos esparsos de magnitude muito
    elevada ($CV$ entre $2.83$ e $5.39$). O estoque inicial $I_0$
    calculado pela Equação~\eqref{eq:I0} cobre os primeiros ciclos,
    o que garante que o diferenciador entre políticas seja
    exclusivamente o custo de manutenção e a frequência de pedidos,
    não o nível de serviço.
  \item Grupo II (6 lojas, $50 < \hat{\mu} \leq 500$):
    estrato de maior relevância prática, com variabilidade entre
    $2.83$ e $4.51$ e volume suficiente para que as políticas se
    diferenciem em todas as cinco métricas avaliadas. É o grupo
    mais representativo das pequenas redes de varejo regional
    do Nordeste brasileiro.
  \item Grupo III (4 lojas, $\hat{\mu} \leq 50$):
    lojas de baixíssimo volume, com até 24 ciclos positivos mas
    demanda acumulada muito baixa. O CV próximo de 3 cria uma barreira
    estrutural: qualquer pedido compatível com a média histórica
    pode ser maior do que a demanda real de vários ciclos futuros,
    tornando a otimização de inventário essencialmente um problema
    de previsão de picos raros.
\end{itemize}

Com o regime de demanda caracterizado, a próxima seção descreve as
doze políticas avaliadas e seus respectivos hiperparâmetros.

% =========================================================
\section{Políticas de Controle de Inventário}
% =========================================================

Esta seção apresenta as doze políticas organizadas em quatro
categorias de crescente sofisticação computacional. A
Tabela~\ref{tab:params} consolida os hiperparâmetros de todas as
abordagens, facilitando reprodutibilidade.
% \subsection{Modelos Clássicos de Inventário}

\subsection{Políticas Clássicas}

O Lote Econômico de Compra (EOQ), introduzido por \citet{harris:1913},
é o modelo analítico fundador do controle de estoques: minimiza a
soma dos custos de ordenação e manutenção sob demanda contínua e
constante, derivando $Q^* = \sqrt{2Dc_o/h}$. Apesar de sua
elegância matemática, a hipótese de estacionariedade o torna
inadequado para séries intermitentes~\citep{silver:2017}.
\citet{hadley:1963} estenderam o modelo para demanda estocástica,
derivando a política $(s,S)$, na qual um pedido de magnitude $S-I_t$
é emitido quando $I_t \leq s$; essa política continua dominante no
varejo regional brasileiro, apesar de igualmente dependente de
estacionariedade. 

A política do Jornaleiro (\textit{Newsvendor}),
formalizada por \citet{cachon:2009}, resolve o problema de período
único como um quantil da distribuição de demanda, equilibrando custo
de excesso e ruptura; sua extensão a horizontes multi-período exige
estimação distribucional que também falha sob intermitência.

\emph{EOQ.} $Q^* = \sqrt{2 D c_o^{fix} / h}$, com ponto de
reposição $s = \hat{\mu} L$~\citep{harris:1913}.

\emph{$(s,S)$.} Pedido de magnitude $S - I_t$ sempre que
$I_t \leq s$, com $s = s_\text{EOQ}$ e
$S = s + Q_\text{EOQ}$~\citep{hadley:1963}.

\emph{Jornaleiro.} $Q^* = \hat{\mu} + z\hat{\sigma}$,
$z = 1.28$ (nível de serviço alvo de $90\%$)~\citep{cachon:2009}.

\subsection{Meta-heurísticas de Otimização}

% \subsection{Algoritmo Genético}

O Algoritmo Genético (GA), proposto por \citet{holland:1975}, é uma
metaheurística evolutiva inspirada na seleção natural: mantém uma
população de soluções candidatas e itera por seleção, cruzamento e
mutação, convergindo para regiões de maior aptidão sem exigir
continuidade nem gradientes. No contexto de inventário, o GA é
empregado para otimização conjunta dos parâmetros $(s, Q)$, tratando
o custo total simulado como função objetivo implícita~\citep{anna:2006}.
\citet{zabraoui:2025} utilizaram o GA precisamente nesse papel:
busca global eficiente de parâmetros estáticos, gerando um ponto de
partida de alta qualidade para o agente de RL subsequente.

% \subsection{Simulated Annealing}

O \textit{Simulated Annealing} (SA), introduzido por
\citet{kirkpatrick:1983}, inspira-se no resfriamento lento de metais
para escapar de ótimos locais: a cada iteração, uma solução vizinha é
aceita deterministicamente se melhor, ou com probabilidade
$\exp(-\Delta C / T_k)$ se pior, onde $T_k$ decresce geometricamente.
A capacidade de escapar de ótimos locais torna o SA atraente para
otimização de parâmetros de inventário em séries de alta variabilidade,
onde a superfície de custo tende a ser multimodal.

% \subsection{Otimização por Enxame de Partículas}

A Otimização por Enxame de Partículas (PSO), proposta por
\citet{kennedy:1995}, coordena um enxame de agentes que exploram
coletivamente o espaço de busca. Cada partícula combina sua melhor
posição histórica (componente cognitivo) com a melhor posição global
(componente social) para atualizar velocidade e posição. O PSO foi
aplicado com sucesso ao controle de estoques por sua convergência
rápida em espaços de parâmetros contínuos~\citep{anna:2006},
apresentando bom desempenho no regime de dimensão moderada, como os
parâmetros $(s, Q)$ individuais por loja deste estudo.

% \subsection{Evolução Diferencial}

A Evolução Diferencial (ED), proposta por \citet{storn:1997}, gera
candidatos por perturbação diferencial de vetores da população. Na
estratégia \textit{rand/1/bin}, um vetor mutante é formado somando ao
vetor base a diferença ponderada de dois outros vetores aleatórios;
o cruzamento binomial combina o mutante com o vetor atual, e o
candidato substitui o pai apenas se melhor. Em comparação com GA e
PSO, a ED tende a convergir mais eficientemente em espaços de baixa
dimensão com função objetivo custosa de avaliar.

Todas as quatro meta-heurísticas buscam $(s^*, Q^*)$ minimizando o custo total acumulado
da Equação~\eqref{eq:obj} ao longo do horizonte de simulação.

\subsection{Agentes de Aprendizado por Reforço}

% \subsection{Rede Q Profunda (DQN)}

A Rede Q Profunda (DQN), proposta por \citet{mnih:2015}, foi o
primeiro método a combinar \textit{Q-learning} com redes neurais em escala,
introduzindo o \textit{experience replay} e a \textit{target network} para
estabilizar o treinamento. \citet{oroojlooyjadid:2022} demonstraram
que a DQN aprende políticas de reposição sem modelar explicitamente
a distribuição de demanda, mas identificaram uma limitação crítica:
o agente exige da ordem de $10^6$ interações para convergência típica,
tornando-o inadequado para horizontes curtos com dados históricos
limitados (como os 38 ciclos deste estudo). Essa limitação é o
principal mecanismo que motiva as arquiteturas híbridas.

% \subsection{Otimização de Política Proximal (PPO)}

A Otimização de Política Proximal (PPO), proposta por
\citet{schulman:2017}, treina uma política estocástica parametrizada
por descida de gradiente com \textit{clipping}: a razão entre a nova e
a antiga probabilidade de ação é truncada no intervalo
$[1-\epsilon_\text{clip}, 1+\epsilon_\text{clip}]$, impedindo
atualizações excessivamente grandes que desestabilizam o aprendizado.
Nativo de ações contínuas e com convergência mais estável que a DQN,
o PPO é, contudo, mais sensível ao ponto de partida do treinamento,
o que o torna particularmente adequado para a fase de refino em
arquiteturas híbridas com inicialização por GA.

% \subsection{SARSA}

O SARSA (\textit{State-Action-Reward-State-Action}), formalizado por
\citet{sutton:2018}, é um algoritmo de RL \textit{on-policy} tabular que
atualiza $Q(s_t, a_t) \leftarrow Q(s_t, a_t) + \alpha[r_t +
\gamma Q(s_{t+1}, a_{t+1}) - Q(s_t, a_t)]$. Por ser \textit{on-policy},
incorpora o custo da própria política de exploração nos valores
aprendidos, tornando-o mais conservador e potencialmente mais
estável que o \textit{Q-learning} em ambientes ruidosos. Sua representação
tabular é computacionalmente leve, mas limita a expressividade em
espaços de estado de alta dimensão; em horizontes curtos com
exploração insuficiente, pode convergir para políticas degeneradas.

Os três agentes aprendem a política diretamente da interação com
o ambiente, sem especificação de $(s, Q)$. O vetor de estado
contém seis componentes normalizados: nível de estoque atual,
demanda esperada no próximo ciclo, pedidos em trânsito, média
móvel e desvio padrão dos últimos seis ciclos, e progresso
relativo do horizonte.

\subsection{Abordagens Híbridas}

A motivação para arquiteturas híbridas decorre da complementaridade
entre meta-heurísticas e agentes de RL: as primeiras realizam busca
global eficiente em parâmetros estáticos; os segundos aprendem
políticas adaptativas estado-dependentes, mas requerem muitas
interações para convergir. \citet{zabraoui:2025} formalizaram esse
princípio na arquitetura GA-DQN, que este trabalho replica e estende
com a variante GA-PPO. A avaliação dessas arquiteturas em regime de
demanda extremamente intermitente é a principal contribuição empírica
do presente estudo: \citet{croston:1972} e \citet{syntetos:2001}
estabeleceram que padrões com CV $> 1.5$ exigem tratamento
metodológico específico, não endereçado por \citet{zabraoui:2025}.

Em ambas as arquiteturas, o \textit{pipeline} opera em três fases
sequenciais: (i)~o GA executa 50 gerações de busca global e identifica
parâmetros estáticos de alta qualidade; (ii)~as soluções de maior
aptidão são usadas para pré-popular um \textit{replay buffer} com
1.000 transições representativas; (iii)~o agente de RL (DQN ou PPO)
treina por 200 episódios a partir desse ponto de partida qualificado,
refinando localmente as decisões a cada ciclo sem precisar explorar
o espaço de estados do zero.

\begin{table}[ht]
\centering
\small
\caption{Hiperparâmetros das doze políticas de controle de
inventário.}
\label{tab:params}
\begin{tabular}{@{}llp{5.5cm}@{}}
\toprule
Política & Categoria & Hiperparâmetros Principais \\
\midrule
EOQ
  & Clássica
  & $Q^* = \sqrt{2Dc_o^{fix}/h}$;\; $s = \hat{\mu}L$ \\
$(s,S)$
  & Clássica
  & $s = s_\text{EOQ}$;\; $S = s + Q_\text{EOQ}$ \\
Jornaleiro
  & Clássica
  & $Q^* = \hat{\mu} + 1.28\hat{\sigma}$;\; SL alvo = $90\%$ \\
\midrule
GA
  & Metaheurística
  & Pop.\ = 100;\; Ger.\ = 50;\; $p_c = 0.7$;\; torneio $k=3$;\; mut.\ gaussiana \\
SA
  & Metaheurística
  & $T_0 = 1000$;\; $\alpha = 0.95$;\; aceitação: $\exp(-\Delta C/T_k)$ \\
PSO
  & Metaheurística
  & $N = 40$ part.;\; 80 iter.;\; $w = 0.7$;\; $c_1 = c_2 = 1.5$ \\
ED
  & Metaheurística
  & \textit{rand/1/bin};\; $F = 0.8$;\; $C_R = 0.9$ \\
\midrule
DQN
  & Reforço
  & 2 camadas $\times$ 64 ReLU;\; 500 ep.;\; $\epsilon$-greedy decrescente;\; \textit{replay} 10.000 \\
PPO
  & Reforço
  & 500 ep.;\; $\epsilon_\text{clip} = 0.2$;\; ações contínuas \\
SARSA
  & Reforço
  & Tabela-$Q$ 5$\times$10;\; $\alpha = 0.1$;\; $\gamma = 0.99$;\; $\epsilon = 0.1$ \\
\midrule
GA-DQN
  & Híbrida
  & GA (50 ger.) $\to$ \textit{buffer} 1.000 trans.\ $\to$ DQN 200 ep. \\
GA-PPO
  & Híbrida
  & GA (50 ger.) $\to$ \textit{buffer} 1.000 trans.\ $\to$ PPO 200 ep. \\
\bottomrule
\end{tabular}
\end{table}

Com as doze políticas e seus hiperparâmetros definidos, a próxima
seção descreve o protocolo de simulação e os critérios de avaliação
comparativa.

% =========================================================
\section{Ambiente de Simulação e Metodologia Experimental}
% =========================================================

Esta seção descreve o ambiente de simulação, a estratégia de
calibração adaptativa que garante comparação justa entre lojas de
volumes muito distintos e os indicadores de desempenho adotados.

\subsection{Ambiente de Simulação}

O ambiente \textit{InventoryEnv}, implementado em Python, simula um
sistema de inventário $(s, Q)$ com prazo de entrega fixo $L = 2$
ciclos. A série de demanda real de cada loja é mantida idêntica para
todas as políticas, de modo que as diferenças de desempenho
reflitam exclusivamente as características de cada política.

\subsection{Calibração Adaptativa de Parâmetros por Loja}

A heterogeneidade entre as lojas ($11.7$ a $1455$ unid./ciclo)
inviabiliza o uso de parâmetros absolutos uniformes. Os parâmetros
de referência são escalonados pelas estatísticas $\hat{\mu}$ e
$\hat{\sigma}$ de cada loja:
\begin{align}
  I_0 &= \max\!\bigl(100,\; 2\hat{\mu} L\bigr), \label{eq:I0}\\
  Q_\text{ref} &= \sqrt{2\hat{\mu} T\, c_o^{fix} / h}, \label{eq:qref}\\
  s_\text{ref} &= \hat{\mu} L + z\hat{\sigma}\sqrt{L}, \label{eq:sref}
\end{align}
com $z = 1.28$ (nível de serviço alvo de $90\%$). Os espaços de
busca das meta-heurísticas e os limites de ação dos agentes de RL
são dimensionados proporcionalmente a $Q_\text{ref}$ e $s_\text{ref}$.

\subsection{Parâmetros de Custo e Indicadores de Desempenho}

Os parâmetros de custo são uniformes para todas as lojas e políticas:
$h = 1.0$, $c_s = 5.0$, $c_o^{fix} = 50.0$ e
$c_o^{var} = 0.5$. Cinco indicadores são computados ao fim dos
38 ciclos:

\begin{itemize}
  \item \emph{CTI} (Custo Total de Inventário):
    $\sum_{t}[h I_t + c_s S_t + c_o^{fix}\mathbf{1}(Q_t>0) + c_o^{var}Q_t]$;
  \item \emph{NS} (Nível de Serviço):
    $\sum_t \min(I_t, D_t) / \sum_t D_t$;
  \item \emph{TR} (Taxa de Ruptura): fração de ciclos com \textit{stockout};
  \item \emph{BE} (\textit{Bullwhip Effect}):
    $\mathrm{Var}(Q_1,\ldots,Q_T)/\mathrm{Var}(D_1,\ldots,D_T)$,
    razão entre a variância dos pedidos e a variância da demanda,
    mede a amplificação da variabilidade ao longo da cadeia;
  \item \emph{FP} (Frequência de Pedidos): fração de ciclos com
    ao menos um pedido emitido.
\end{itemize}

Com o protocolo experimental definido, a próxima seção apresenta
e discute os resultados comparativos das doze políticas.

% =========================================================
\section{Resultados e Análise}
% =========================================================

Esta seção apresenta a comparação global das doze políticas, a
análise estratificada por grupo de demanda, o comportamento do efeito
\textit{bullwhip} e a identificação da política de melhor equilíbrio,
com discussão dos mecanismos subjacentes e comparação direta com
\citet{zabraoui:2025}.

\subsection{Comparação Global entre Políticas}

A Tabela~\ref{tab:global} apresenta as médias dos cinco indicadores
sobre as quinze lojas, ordenadas por CTI crescente. Os melhores
valores em cada coluna estão destacados em negrito.

\begin{table}[ht]
\centering
\small
\caption{Médias dos indicadores de desempenho por política (15 lojas).
Melhores valores por coluna em \textbf{negrito}. $^\dagger$SARSA e PPO
apresentam degeneração ($FP = 0$ e $FP = 0.98$, respectivamente).}
\label{tab:global}
\begin{tabular}{@{}lrrrrr@{}}
\toprule
Política &
\makecell[r]{CTI\\(R\$)} &
\makecell[r]{NS} &
\makecell[r]{TR} &
\makecell[r]{BE} &
\makecell[r]{FP} \\
\midrule
SARSA$^\dagger$            & \textbf{16.139} & 0{,}785 & 0{,}200 & \textbf{0{,}00}  & \textbf{0{,}00} \\
Jornaleiro (\textit{Newsvendor})    & 16.402 & 0{,}794 & 0{,}111 & 0{,}003 & 0{,}27 \\
Rede Q Profunda (DQN)      & 17.974 & 0{,}806 & 0{,}078 & 4{,}1   & 0{,}37 \\
GA-DQN (Híbrido)           & 18.502 & 0{,}821 & 0{,}050 & 5{,}0   & 0{,}31 \\
Enxame de Partículas (PSO) & 19.100 & 0{,}819 & 0{,}056 & 6{,}2   & 0{,}31 \\
GA-PPO (Híbrido)           & 19.801 & 0{,}833 & 0{,}050 & 6{,}6   & 0{,}35 \\
Algoritmo Genético (GA)    & 20.015 & 0{,}825 & 0{,}061 & 19{,}7  & 0{,}30 \\
Evolução Diferencial (ED)  & 20.252 & 0{,}813 & 0{,}056 & 31{,}0  & 0{,}27 \\
\textit{Sim. Annealing} (SA) & 20.967 & 0{,}820 & 0{,}056 & 36{,}4  & 0{,}27 \\
Lote Econômico (EOQ)       & 31.034 & 0{,}840 & 0{,}050 & 23{,}0  & 0{,}20 \\
Política PPO$^\dagger$     & 34.637 & \textbf{0{,}848} & \textbf{0{,}039} & 15{,}9 & 0{,}98 \\
Política $(s,S)$           & 35.589 & 0{,}843 & 0{,}050 & 119{,}4 & 0{,}15 \\
\bottomrule
\end{tabular}
\end{table}

Os dois extremos revelam os riscos da otimização por única métrica:
o SARSA converge à estratégia de não emitir pedidos (FP\,=\,0),
obtendo o menor CTI ao preço da maior taxa de ruptura ($0.200$); o
PPO emite pedidos em $98\%$ dos ciclos, atingindo o maior NS mas
incorrendo em CTI inviável de R\$\,34.637. Entre as políticas
práticas (sem degeneração), o GA-DQN apresenta o melhor equilíbrio
global, com CTI apenas $13\%$ acima do Jornaleiro e TR $55\%$ menor.
A política $(s,S)$, ainda amplamente adotada no varejo regional,
tem CTI $92\%$ superior ao GA-DQN e efeito \textit{bullwhip} de 119
vezes, configurando uma escolha indefensável nesse regime.

\subsection{Análise Estratificada por Grupo de Demanda}

A estratificação por grupo evidencia que não existe política
universalmente dominante, resultado de maior relevância prática
para gestores de \textit{supply chain}.

Grupo I ($\hat{\mu} > 500$, cinco lojas): estoque inicial
adequado garante NS\,=\,$1.00$ para quase todas as políticas; o
diferenciador é o custo. O EOQ gasta em média R\$\,85.956 por loja
contra R\$\,45.739 das políticas de menor frequência de pedidos,
redução de $47\%$ apenas pela troca de heurística.

Grupo II ($50 < \hat{\mu} \leq 500$, seis lojas): estrato
mais discriminante e de maior relevância prática. Nas médias do
Grupo~II, SA, PSO e ED atingem NS entre $0.89$ e $0.93$ com CTI
entre R\$\,15.465 e R\$\,20.252. O GA-DQN destaca-se com
NS\,=\,$0.933$ e CTI\,=\,R\$\,10.854, a melhor combinação
custo-serviço do estudo, atribuível ao \textit{replay buffer} de
alta qualidade gerado pelo GA~\citep{oroojlooyjadid:2022}.

Grupo III ($\hat{\mu} \leq 50$, quatro lojas): três lojas
não superam NS de $0.42$ com nenhuma política. O CV próximo de 3
cria uma barreira estrutural em que pedidos compatíveis com a média
não cobrem picos e pedidos elevados aumentam o custo sem ganho de
serviço. Nesse estrato, o Jornaleiro é a política mais econômica.

\subsection{Efeito de Amplificação da Demanda (\textit{Bullwhip})}

O efeito \textit{bullwhip}~\citep{lee:1997} mede quanto a
variabilidade dos pedidos supera a da demanda; valores elevados
implicam oscilações desnecessárias e custos logísticos crescentes a
montante. Os resultados revelam amplitude expressiva: excluindo o SARSA
(degenerado, BE\,=\,0), as políticas práticas variam de $0.003$
(Jornaleiro) a $119.4$ em $(s,S)$, razão superior a 40.000 vezes
e equivalente a quase cinco ordens de grandeza. Entre as políticas
práticas, DQN ($4.1$) e GA-DQN ($5.0$) apresentam menor
amplificação por aprenderem a antecipar picos de demanda; PSO ($6.2$)
é a melhor metaheurística nessa dimensão, enquanto SA ($36.4$) e ED
($31.0$) ficam elevados por otimizarem $(s,Q)$ pela média histórica.
Este indicador, presente em \citet{zabraoui:2025}
apenas para políticas selecionadas, é aqui sistematizado como critério
primário para todas as doze políticas, tornando explícito o impacto
de cada abordagem sobre a estabilidade de toda a cadeia.

A Figura~\ref{fig:barras} sintetiza as três métricas primárias.

\begin{figure}[ht]
  \centering
  \includegraphics[width=\linewidth]{figuras/pb_barras.png}
  \caption{Desempenho comparativo das doze políticas segundo CTI
  médio, NS médio e BE médio (barras ordenadas do melhor para o pior
  em cada indicador).}
  \label{fig:barras}
\end{figure}

\subsection{Política Recomendada e Comparação com Zabraoui (2025)}

Considerando os cinco indicadores conjuntamente, o GA-DQN emerge
como a política de melhor equilíbrio: CTI de R\$\,18.502 (apenas
$13\%$ acima do Jornaleiro, que tem TR mais que o dobro), NS de $0.821$
equiparável ao das meta-heurísticas isoladas ($0.813$--$0.825$) com CTI
significativamente menor,
BE de $5.0$ e TR de $0.050$ equivalente à $(s,S)$ porém com custo
$92\%$ menor. O mecanismo subjacente é a complementaridade entre as
abordagens: o GA identifica $(s,Q)$ de qualidade com apenas
$100 \times 50 = 5.000$ avaliações, suprindo a insuficiência do
DQN em horizontes curtos (da ordem de $10^6$ passos para
convergência típica~\citep{mnih:2015}).

Esses resultados confirmam, em contexto de $CV$ entre $2.5$ e $5.4$,
a vantagem do GA-DQN reportada por \citet{zabraoui:2025} para o
Walmart~M5 (CV típico $< 1$), demonstrando que a arquitetura é
robusta entre regimes de demanda e ampliando significativamente
seu escopo de validade.

% =========================================================
\section{Conclusão}
% =========================================================

Esta seção sintetiza as recomendações práticas, discute o impacto
econômico e de mercado e aponta as limitações e direções futuras
deste trabalho.

A comparação sistemática de doze políticas em dados reais de varejo
brasileiro com $CV$ entre $2.5$ e $5.4$ demonstrou que a arquitetura
híbrida GA-DQN oferece o melhor equilíbrio entre custo total,
nível de serviço e estabilidade logística nesse regime, superando
tanto as heurísticas clássicas quanto os agentes de aprendizado por
reforço isolados. A política $(s,S)$, amplamente adotada no varejo
regional brasileiro, mostrou-se indefensável: é dominada em custo e
estabilidade por todas as alternativas práticas avaliadas, com efeito
\textit{bullwhip} 24 vezes maior que o GA-DQN e CTI $92\%$ superior.

\subsection{Recomendações Operacionais}

Com base na análise estratificada, três recomendações práticas
emergem diretamente dos resultados:

\begin{itemize}
  \item Adotar GA-DQN como política padrão para o Grupo II
    (lojas de volume médio): NS\,=\,$0.933$ e
    CTI\,=\,R\$\,10.854 justificam o custo computacional adicional
    frente às heurísticas clássicas. O treinamento é realizado uma
    única vez por loja e produto.
  \item Substituir imediatamente a política $(s,S)$: os resultados
    da Seção~6 demonstram que ela é dominada em custo e estabilidade
    logística por todas as alternativas práticas, sendo indefensável
    em qualquer estrato de demanda extremamente intermitente. A
    migração direta para o Jornaleiro, sem custo computacional
    adicional, já elimina o principal problema de amplificação da
    cadeia.
  \item Utilizar o Jornaleiro no Grupo~III (baixo volume):
    sua simplicidade e baixo CTI superam abordagens sofisticadas
    nesse estrato, onde a barreira estrutural de NS limita qualquer
    política independentemente de sua complexidade.
\end{itemize}

\subsection{Contribuição para o Mercado e Importância da Pesquisa}

Esta pesquisa tem relevância direta para três segmentos de mercado.
Primeiro, para redes de varejo regional com produtos de baixo
giro: a substituição da política $(s,S)$ pelo GA-DQN pode representar
redução de custo de inventário da ordem de $48\%$ sem perda de nível
de serviço, permitindo realocar capital antes imobilizado em estoque
para outras prioridades operacionais. Em termos absolutos, a diferença
de CTI entre $(s,S)$ e GA-DQN é de R\$\,17.087 por loja por período;
extrapolada a uma rede de 100 lojas com 50 produtos de comportamento
similar, representa uma economia potencial da ordem de
R\$\,85 milhões por período de 38 ciclos. Segundo, para
distribuidores e operadores logísticos: o controle do efeito
\textit{bullwhip} (reduzido de $119$ para $5$) estabiliza a cadeia a
montante, diminuindo custos de transporte e de capacidade de
armazenagem nos centros de distribuição. Terceiro, para
pequenas e médias empresas sem equipe de \textit{data science}:
a arquitetura GA-DQN requer apenas dados históricos de venda (já
disponíveis em qualquer sistema ERP, \textit{Enterprise Resource Planning})
e pode ser operacionalizada
como rotina periódica automatizada.

Do ponto de vista científico, o estudo oferece à comunidade de
Pesquisa Operacional uma base de comparação reproduzível em regime
de extrema intermitência, preenchendo a lacuna deixada pelos
estudos baseados em Walmart~M5. A metodologia é replicável para
qualquer produto com CV $> 1.5$ em qualquer segmento de varejo.

\subsection{Limitações e Direções Futuras}

O estudo adota simulação determinística com série de demanda estática,
sem réplicas independentes. Como direções futuras, propõe-se:
avaliação \textit{out-of-sample} com janela deslizante para validar a
generalização; extensão a múltiplos produtos com restrições de
capacidade de armazenagem; incorporação de prazo de entrega
estocástico; e integração com modelos de previsão de demanda (Croston,
\textit{Long Short-Term Memory} (LSTM)) para alimentar as decisões com estimativas prospectivas.

\bibliographystyle{sbpo}
\bibliography{exemplo-latex}

\end{document}
